% A. Introducción al problema (Guía apdo. 3.A)
% Extensión orientativa: 3-4 páginas
\chapter{Introducción}
\label{cap:introduccion}

La monitorización de activos críticos resulta fundamental para garantizar la operatividad
y prevenir averías inesperadas en entornos que abarcan desde el sector servicios hasta el
industrial. Los sistemas electromecánicos de refrigeración constituyen un caso
representativo: operan de forma continua, su fallo implica la pérdida del producto
conservado y su degradación rara vez es instantánea, sino que viene precedida por cambios
progresivos en su comportamiento físico.

% TODO: cuantificar el problema con fuentes citables (coste de la parada no planificada,
% porcentaje de fallos en compresores, ahorro atribuido al mantenimiento predictivo).
% Añadir las entradas correspondientes a referencias.bib en formato IEEE.

\section{Motivación}

El mantenimiento correctivo actúa una vez que el fallo ya se ha producido, mientras que el
mantenimiento preventivo basado en calendario sustituye componentes que a menudo conservan
vida útil. El mantenimiento predictivo se apoya en la medida continua de variables físicas
para estimar el estado real del activo y anticipar la intervención únicamente cuando el
comportamiento se desvía del nominal.

La adopción de este paradigma se ve limitada, en la práctica, por el coste de la
infraestructura de captura y por el volumen de datos que genera un muestreo de señales
mecánicas y acústicas. Enviar de forma continua señales de vibración a la nube resulta
costoso en ancho de banda y en almacenamiento, e introduce una dependencia de conectividad
que un activo industrial no siempre puede asumir.

\section{Planteamiento del problema}

Este trabajo aborda el diseño e implementación de un sistema de monitorización modular que
desplaza la lógica de diagnóstico hacia el extremo de la red (\textit{edge}). El nodo de
medida captura variables físicas del activo --vibración, temperatura y firma acústica-- y
el objetivo último es que la detección de anomalías se resuelva en el propio nodo, sin
depender de un flujo masivo de datos hacia la nube.

El activo seleccionado como escenario de prueba es el compresor de un sistema de
refrigeración, por su disponibilidad, su funcionamiento cíclico bien definido y la
posibilidad de inducir fallos controlados sin riesgo.

\section{Estructura del documento}

El capítulo~\ref{cap:estado-arte} analiza las necesidades del problema y el estado del arte
en mantenimiento predictivo e inteligencia en el borde. El capítulo~\ref{cap:objetivos}
detalla los objetivos perseguidos. El capítulo~\ref{cap:metodologia} describe la
metodología seguida de forma sucesiva. El capítulo~\ref{cap:diseno} justifica las
decisiones técnicas adoptadas, incluidos los estándares empleados. El
capítulo~\ref{cap:resultados} presenta los resultados obtenidos y el
capítulo~\ref{cap:conclusiones} recoge la discusión, las conclusiones y las líneas futuras.
Los anexos técnicos completan la información de detalle del sistema construido.
